\scanpage{91}
% archive sheet number: 90
% block text | lines 1-2
Or, reproduisant un raisonnement déjà fait, on en conclut que $H^1(\Gamma,A^*)=\{e\}$.

% block text | line 3
On est donc ramené à étudier
% block display | line 4
\[
H^1(\Gamma,A^*)
\]
% block text | lines 5-9
quand $O$ est un \uline{anneau local}, $\Gamma$ opère \uline{\unsure{trivialt|trivlt}} sur $O/m$, et $A$ est une $(O$-$\Gamma)$-algèbre. \uline{Supposons} \inserted{\uline{de plus que $\Gamma$ opère trivialement sur $A/mA$.}} \inserted{et soit} Soit $p$ la caractéristique de $k=O/m$, $\Lambda$ un sous-groupe de Sylow du type $p$ de $\Gamma$ ($\Lambda=\{e\}$ si la caractéristique est nulle).
% block margin note | left margin beside lines 5-17
% slanted
\margin{Nous supposerons plus précisément que $A^\Gamma\longrightarrow A/mA$ est \uuline{surjectif} [ce qui est p.ex. le cas si $A=M_n(\underline{O})$, $\Gamma$ opérant \unsure{trivialt|trivialement} sur $\underline{O}/m$, en effet, on sait qu'alors $\underline{O}^\Gamma\longrightarrow\underline{O}/m$ est surjectif\inserted{??}.....]}
% block text | line 10
Je dis qu'alors l'application
% block diagram | line 11
\diagram{\begin{tikzcd}[ampersand replacement=\&]
H^1(\Gamma,A^*)\arrow[r]\&H^1(\Gamma,A/mA^*)\times H^1(\Lambda,A^*)
\end{tikzcd}}
% block text | lines 12-15
est \deleted{bijective} injective (en particulier, si l'ordre de $G$ est premier à la car. $p$, en particulier si $k$ est de car. $0$, alors $H^1(\Gamma,A^*)\mathbin{-}H^1(\Gamma,A/mA^*)$ est injectif)

% block text | lines 16-22
Soient $f$ et $g$ \deleted{\unsure{enfin}} deux cocycles donnant les $=^{\mathrm{es}}$ images dans $H^1(\Gamma,A/mA^*)$ et $H^1(\Lambda,A^*)$. Pour montrer qu'ils sont $\sim$, on peut \unsure{encore|comme} supposer $g\unsure{(\gamma)|\gamma|(\sigma)}=\unsure{e|1}$ pour $\forall\gamma$. \deleted{\unsure{Soit le}} \inserted{Il suffit de} prouver alors \deleted{\unsure{$f\sim0$}} que $f$ est $\sim0$ dans \deleted{\unsure{$H^1(\Gamma,$|$H^1$},} $A/m^n A^*$ pour tout $n$. On \unsure{récurre} sur $n$, supposant que c'est fait pour un $n\geq1$.
